\begin{table}[htbp]
\centering
\caption{Effect of USPS Establishment Losses on Presidential Election Turnout}
\label{tab:main}
\begin{tabular}{lccc}
\hline\hline
 & (1) & (2) & (3) \\
 & CS-DiD & TWFE & TWFE Intensity \\
\hline
Treatment & $-0.0141$ & $-0.0254$ & $-0.0232$ \\
 & $(0.0058)$ & $(0.0114)$ & $(0.0091)$ \\
 & $[-0.0255, -0.0028]$ & $[-0.0477, -0.0030]$ & $[-0.0411, -0.0053]$ \\
\hline
Observations & 12,537 & 12,537 & 12,537 \\
Counties & 1791 & 1791 & 1791 \\
Treated counties & 359 & 359 & 359 \\
Estimator & Callaway-Sant'Anna & TWFE & TWFE \\
Control group & Never-treated & --- & --- \\
PT pre-test $p$-value & $0.0035$ & --- & --- \\
Clustering & Analytical & State & State \\
\hline\hline
\end{tabular}
\begin{tablenotes}
\small
\item \textit{Notes:} Dependent variable is log total presidential votes at the county level. Panel spans seven presidential elections (2000--2024). Column (1) reports the overall ATT from the Callaway and Sant'Anna (2021) estimator using doubly robust estimation with never-treated counties as the control group. Columns (2)--(3) report two-way fixed effects estimates. Treatment in (1)--(2) is binary (county lost $\geq 1$ USPS establishment, 2015--2018). Treatment in (3) is continuous (number of establishments lost). Standard errors in parentheses; 95\% confidence intervals in brackets. Column (2)--(3) SEs clustered at the state level (29 states). The parallel trends pre-test in column (1) rejects at the 1\% level, suggesting a pre-existing convergence trend.
\end{tablenotes}
\end{table}
